% ==============================================================================
% Figure 1: LCESA Framework Overview
% Three-layer instantiation: Abstract Geometry -> Five-Layer Model -> Transformer
%
% Usage: % ==============================================================================
% Figure 1: LCESA Framework Overview
% Three-layer instantiation: Abstract Geometry -> Five-Layer Model -> Transformer
%
% Usage: % ==============================================================================
% Figure 1: LCESA Framework Overview
% Three-layer instantiation: Abstract Geometry -> Five-Layer Model -> Transformer
%
% Usage: % ==============================================================================
% Figure 1: LCESA Framework Overview
% Three-layer instantiation: Abstract Geometry -> Five-Layer Model -> Transformer
%
% Usage: \input{figures/fig1_lcesa_framework.tex}
% Requires: \usepackage{tikz}
%           \usetikzlibrary{arrows.meta, positioning, shapes.geometric,
%                           decorations.pathreplacing, calc, fit}
% ==============================================================================

\begin{figure}[t]
\centering
\begin{tikzpicture}[
    % --- global styles ---
    >=Stealth,
    node distance=0.6cm and 0.4cm,
    every node/.style={font=\small},
    % --- layer box styles ---
    layerbox/.style={
        draw=#1, fill=#1!6, rounded corners=4pt,
        inner sep=8pt, minimum width=12.8cm
    },
    % --- component box styles ---
    abstractbox/.style={
        draw=abstractclr, fill=abstractclr!10, rounded corners=3pt,
        inner sep=5pt, minimum height=0.9cm, font=\small
    },
    layercomp/.style={
        draw=layerclr, fill=layerclr!10, rounded corners=3pt,
        inner sep=4pt, minimum height=0.85cm, minimum width=1.9cm,
        font=\small, align=center
    },
    instbox/.style={
        draw=instclr, fill=instclr!10, rounded corners=3pt,
        inner sep=4pt, minimum height=0.85cm, minimum width=2.1cm,
        font=\small, align=center
    },
    % --- arrow styles ---
    maparrow/.style={
        ->, line width=0.9pt, draw=#1, >=Stealth
    },
    maplabel/.style={
        font=\footnotesize, text=#1, fill=white, inner sep=1.5pt
    },
]

% === Colors ===
\definecolor{abstractclr}{RGB}{41, 98, 165}   % blue for abstract geometry
\definecolor{layerclr}{RGB}{34, 139, 87}      % green for five-layer model
\definecolor{instclr}{RGB}{180, 95, 25}       % orange for transformer instantiation
\definecolor{arrowclr}{RGB}{100, 100, 100}

% ====================================================================
% LAYER 1: Abstract Geometric Framework
% ====================================================================
\node[layerbox=abstractclr, label={[font=\small\bfseries, abstractclr]above:Abstract Geometric Framework}]
    (layer1) at (0, 7.4) {};

% Components inside layer 1
\node[abstractbox, minimum width=2.4cm] (eventspace)
    at (-4.6, 7.4) {$\mathcal{X} = (V, E, T)$};
\node[abstractbox, minimum width=1.6cm] (fiber)
    at (-1.6, 7.4) {$E_x$};
\node[abstractbox, minimum width=2.0cm] (connection)
    at (1.1, 7.4) {$A_{\Psi}$};
\node[abstractbox, minimum width=1.6cm] (curvature)
    at (3.6, 7.4) {$F$};
\node[abstractbox, minimum width=2.0cm] (gauge)
    at (5.8, 7.4) {$G = \prod_k O(d_k)$};

% Sub-labels beneath each component
\node[font=\scriptsize, text=abstractclr] at (-4.6, 6.85) {event space};
\node[font=\scriptsize, text=abstractclr] at (-1.6, 6.85) {fiber};
\node[font=\scriptsize, text=abstractclr] at (1.1, 6.85) {connection};
\node[font=\scriptsize, text=abstractclr] at (3.6, 6.85) {curvature};
\node[font=\scriptsize, text=abstractclr] at (5.8, 6.85) {gauge group};

% Arrows inside layer 1 (horizontal flow)
\draw[->, abstractclr, line width=0.6pt] (eventspace.east) -- (fiber.west);
\draw[->, abstractclr, line width=0.6pt] (fiber.east) -- (connection.west);
\draw[->, abstractclr, line width=0.6pt] (connection.east) -- (curvature.west);
\draw[->, abstractclr, line width=0.6pt] (curvature.east) -- (gauge.west);

% ====================================================================
% LAYER 2: Five-Layer Standpoint Model
% ====================================================================
\node[layerbox=layerclr, label={[font=\small\bfseries, layerclr]above:Five-Layer Standpoint Model $\Psi_t$}]
    (layer2) at (0, 4.8) {};

% Five standpoint layer components
\node[layercomp, minimum width=2.1cm] (psi_min)
    at (-4.7, 4.8) {$\psi_{\min}$\\[-2pt]{\scriptsize minimal self}};
\node[layercomp, minimum width=2.1cm] (psi_nar)
    at (-2.3, 4.8) {$\psi_{\mathrm{nar}}$\\[-2pt]{\scriptsize narrative}};
\node[layercomp, minimum width=2.1cm] (psi_soc)
    at (0.1, 4.8) {$\psi_{\mathrm{soc}}$\\[-2pt]{\scriptsize social}};
\node[layercomp, minimum width=2.1cm] (psi_mor)
    at (2.5, 4.8) {$\psi_{\mathrm{mor}}$\\[-2pt]{\scriptsize moral}};
\node[layercomp, minimum width=2.1cm] (psi_pos)
    at (4.9, 4.8) {$\psi_{\mathrm{pos}}$\\[-2pt]{\scriptsize positional}};

% Orthogonal decomposition label
\node[font=\scriptsize, text=layerclr] at (0, 4.2)
    {$\Psi_t = \psi_{\mathrm{min}} \oplus \psi_{\mathrm{nar}} \oplus \psi_{\mathrm{soc}} \oplus \psi_{\mathrm{mor}} \oplus \psi_{\mathrm{pos}}$};

% ====================================================================
% LAYER 3: Transformer Instantiation
% ====================================================================
\node[layerbox=instclr, label={[font=\small\bfseries, instclr]above:Transformer Instantiation}]
    (layer3) at (0, 2.0) {};

% Components inside layer 3
\node[instbox, minimum width=2.6cm] (turns)
    at (-4.3, 2.0) {Turns $x_i$\\[-2pt]{\scriptsize (events)}};
\node[instbox, minimum width=2.6cm] (subspaces)
    at (-1.3, 2.0) {Head subspaces\\[-2pt]{\scriptsize $W_k = \operatorname{span}(V_h)$}};
\node[instbox, minimum width=2.6cm] (attn)
    at (1.7, 2.0) {Attention $\alpha$\\[-2pt]{\scriptsize (transport $U_{ij}$)}};
\node[instbox, minimum width=2.6cm] (pathdep)
    at (4.7, 2.0) {Path-dependence\\[-2pt]{\scriptsize (curvature $F_{ijk}$)}};

% Sub-labels
\node[font=\scriptsize, text=instclr] at (-4.3, 1.35) {$V = \{x_1, \ldots, x_n\}$};
\node[font=\scriptsize, text=instclr] at (-1.3, 1.35) {$\mathbb{R}^d = \bigoplus_k W_k$};
\node[font=\scriptsize, text=instclr] at (1.7, 1.35) {$U_{ij} = \sum_k \bar\alpha_{ji}^{(k)} P_{W_k}$};
\node[font=\scriptsize, text=instclr] at (4.7, 1.35) {$F_{ijk} = U_{ij} U_{jk} (U_{ik}^{\exp})^{-1}$};

% Arrows inside layer 3
\draw[->, instclr, line width=0.6pt] (turns.east) -- (subspaces.west);
\draw[->, instclr, line width=0.6pt] (subspaces.east) -- (attn.west);
\draw[->, instclr, line width=0.6pt] (attn.east) -- (pathdep.west);

% ====================================================================
% INTER-LAYER MAPPING ARROWS
% ====================================================================

% Layer 1 -> Layer 2 (left side)
\draw[maparrow=arrowclr, dashed] (-3.0, 6.55) -- (-3.0, 5.45)
    node[midway, right=3pt, maplabel=arrowclr]
    {instantiate};
\node[font=\scriptsize, text=arrowclr, anchor=west] at (-2.8, 5.95)
    {$\Psi_{x_i} \in \bigoplus_k W_k$};

% Layer 2 -> Layer 3 (right side)
\draw[maparrow=arrowclr, dashed] (3.0, 4.15) -- (3.0, 3.0)
    node[midway, right=3pt, maplabel=arrowclr]
    {compute};
\node[font=\scriptsize, text=arrowclr, anchor=west] at (3.2, 3.55)
    {attention heads};

% Layer 1 -> Layer 3 (curved arrow, right side, bypassing layer 2)
\draw[maparrow=arrowclr, dotted, line width=0.8pt]
    (5.6, 6.55) to[out=-70, in=70]
    node[right=2pt, midway, font=\scriptsize, text=arrowclr]
    {direct instantiation}
    (5.6, 3.0);

% ====================================================================
% BRACE: "Standpoint Bundle" annotation on the left
% ====================================================================
\draw[decorate, decoration={brace, amplitude=6pt, mirror}, line width=0.7pt, arrowclr]
    (-6.6, 2.8) -- (-6.6, 6.2)
    node[midway, left=8pt, font=\footnotesize\bfseries, text=arrowclr, rotate=90, anchor=south]
    {Standpoint Bundle};

% ====================================================================
% RIGHT ANNOTATION: Key identity
% ====================================================================
\node[draw=arrowclr, fill=white, rounded corners=2pt,
      inner sep=5pt, font=\footnotesize, align=center, text width=3.4cm]
    at (7.2, 4.8)
    {\textbf{Key Identity:}\\[2pt]
     Attention $=$ Parallel\\Transport on $(E, X, \pi)$};

\end{tikzpicture}
\caption{%
    \textbf{LCESA framework overview.}
    The abstract geometric framework (top) defines the event space~$\mathcal{X}$, fiber~$E_x$,
    connection~$A_\Psi$, and curvature~$F$ on the standpoint bundle.
    The five-layer standpoint model (middle) instantiates the fiber as a direct sum of
    psychologically grounded subspaces~$\psi_k$.
    The transformer instantiation (bottom) maps each geometric object to a concrete
    architectural component: events to turns, fibers to head value subspaces~$W_k$,
    connection to attention weights~$\alpha$, and curvature to path-dependence of
    composite transport.
    Dashed arrows indicate the instantiation mapping between layers.
}
\label{fig:lcesa-framework}
\end{figure}

% Requires: \usepackage{tikz}
%           \usetikzlibrary{arrows.meta, positioning, shapes.geometric,
%                           decorations.pathreplacing, calc, fit}
% ==============================================================================

\begin{figure}[t]
\centering
\begin{tikzpicture}[
    % --- global styles ---
    >=Stealth,
    node distance=0.6cm and 0.4cm,
    every node/.style={font=\small},
    % --- layer box styles ---
    layerbox/.style={
        draw=#1, fill=#1!6, rounded corners=4pt,
        inner sep=8pt, minimum width=12.8cm
    },
    % --- component box styles ---
    abstractbox/.style={
        draw=abstractclr, fill=abstractclr!10, rounded corners=3pt,
        inner sep=5pt, minimum height=0.9cm, font=\small
    },
    layercomp/.style={
        draw=layerclr, fill=layerclr!10, rounded corners=3pt,
        inner sep=4pt, minimum height=0.85cm, minimum width=1.9cm,
        font=\small, align=center
    },
    instbox/.style={
        draw=instclr, fill=instclr!10, rounded corners=3pt,
        inner sep=4pt, minimum height=0.85cm, minimum width=2.1cm,
        font=\small, align=center
    },
    % --- arrow styles ---
    maparrow/.style={
        ->, line width=0.9pt, draw=#1, >=Stealth
    },
    maplabel/.style={
        font=\footnotesize, text=#1, fill=white, inner sep=1.5pt
    },
]

% === Colors ===
\definecolor{abstractclr}{RGB}{41, 98, 165}   % blue for abstract geometry
\definecolor{layerclr}{RGB}{34, 139, 87}      % green for five-layer model
\definecolor{instclr}{RGB}{180, 95, 25}       % orange for transformer instantiation
\definecolor{arrowclr}{RGB}{100, 100, 100}

% ====================================================================
% LAYER 1: Abstract Geometric Framework
% ====================================================================
\node[layerbox=abstractclr, label={[font=\small\bfseries, abstractclr]above:Abstract Geometric Framework}]
    (layer1) at (0, 7.4) {};

% Components inside layer 1
\node[abstractbox, minimum width=2.4cm] (eventspace)
    at (-4.6, 7.4) {$\mathcal{X} = (V, E, T)$};
\node[abstractbox, minimum width=1.6cm] (fiber)
    at (-1.6, 7.4) {$E_x$};
\node[abstractbox, minimum width=2.0cm] (connection)
    at (1.1, 7.4) {$A_{\Psi}$};
\node[abstractbox, minimum width=1.6cm] (curvature)
    at (3.6, 7.4) {$F$};
\node[abstractbox, minimum width=2.0cm] (gauge)
    at (5.8, 7.4) {$G = \prod_k O(d_k)$};

% Sub-labels beneath each component
\node[font=\scriptsize, text=abstractclr] at (-4.6, 6.85) {event space};
\node[font=\scriptsize, text=abstractclr] at (-1.6, 6.85) {fiber};
\node[font=\scriptsize, text=abstractclr] at (1.1, 6.85) {connection};
\node[font=\scriptsize, text=abstractclr] at (3.6, 6.85) {curvature};
\node[font=\scriptsize, text=abstractclr] at (5.8, 6.85) {gauge group};

% Arrows inside layer 1 (horizontal flow)
\draw[->, abstractclr, line width=0.6pt] (eventspace.east) -- (fiber.west);
\draw[->, abstractclr, line width=0.6pt] (fiber.east) -- (connection.west);
\draw[->, abstractclr, line width=0.6pt] (connection.east) -- (curvature.west);
\draw[->, abstractclr, line width=0.6pt] (curvature.east) -- (gauge.west);

% ====================================================================
% LAYER 2: Five-Layer Standpoint Model
% ====================================================================
\node[layerbox=layerclr, label={[font=\small\bfseries, layerclr]above:Five-Layer Standpoint Model $\Psi_t$}]
    (layer2) at (0, 4.8) {};

% Five standpoint layer components
\node[layercomp, minimum width=2.1cm] (psi_min)
    at (-4.7, 4.8) {$\psi_{\min}$\\[-2pt]{\scriptsize minimal self}};
\node[layercomp, minimum width=2.1cm] (psi_nar)
    at (-2.3, 4.8) {$\psi_{\mathrm{nar}}$\\[-2pt]{\scriptsize narrative}};
\node[layercomp, minimum width=2.1cm] (psi_soc)
    at (0.1, 4.8) {$\psi_{\mathrm{soc}}$\\[-2pt]{\scriptsize social}};
\node[layercomp, minimum width=2.1cm] (psi_mor)
    at (2.5, 4.8) {$\psi_{\mathrm{mor}}$\\[-2pt]{\scriptsize moral}};
\node[layercomp, minimum width=2.1cm] (psi_pos)
    at (4.9, 4.8) {$\psi_{\mathrm{pos}}$\\[-2pt]{\scriptsize positional}};

% Orthogonal decomposition label
\node[font=\scriptsize, text=layerclr] at (0, 4.2)
    {$\Psi_t = \psi_{\mathrm{min}} \oplus \psi_{\mathrm{nar}} \oplus \psi_{\mathrm{soc}} \oplus \psi_{\mathrm{mor}} \oplus \psi_{\mathrm{pos}}$};

% ====================================================================
% LAYER 3: Transformer Instantiation
% ====================================================================
\node[layerbox=instclr, label={[font=\small\bfseries, instclr]above:Transformer Instantiation}]
    (layer3) at (0, 2.0) {};

% Components inside layer 3
\node[instbox, minimum width=2.6cm] (turns)
    at (-4.3, 2.0) {Turns $x_i$\\[-2pt]{\scriptsize (events)}};
\node[instbox, minimum width=2.6cm] (subspaces)
    at (-1.3, 2.0) {Head subspaces\\[-2pt]{\scriptsize $W_k = \operatorname{span}(V_h)$}};
\node[instbox, minimum width=2.6cm] (attn)
    at (1.7, 2.0) {Attention $\alpha$\\[-2pt]{\scriptsize (transport $U_{ij}$)}};
\node[instbox, minimum width=2.6cm] (pathdep)
    at (4.7, 2.0) {Path-dependence\\[-2pt]{\scriptsize (curvature $F_{ijk}$)}};

% Sub-labels
\node[font=\scriptsize, text=instclr] at (-4.3, 1.35) {$V = \{x_1, \ldots, x_n\}$};
\node[font=\scriptsize, text=instclr] at (-1.3, 1.35) {$\mathbb{R}^d = \bigoplus_k W_k$};
\node[font=\scriptsize, text=instclr] at (1.7, 1.35) {$U_{ij} = \sum_k \bar\alpha_{ji}^{(k)} P_{W_k}$};
\node[font=\scriptsize, text=instclr] at (4.7, 1.35) {$F_{ijk} = U_{ij} U_{jk} (U_{ik}^{\exp})^{-1}$};

% Arrows inside layer 3
\draw[->, instclr, line width=0.6pt] (turns.east) -- (subspaces.west);
\draw[->, instclr, line width=0.6pt] (subspaces.east) -- (attn.west);
\draw[->, instclr, line width=0.6pt] (attn.east) -- (pathdep.west);

% ====================================================================
% INTER-LAYER MAPPING ARROWS
% ====================================================================

% Layer 1 -> Layer 2 (left side)
\draw[maparrow=arrowclr, dashed] (-3.0, 6.55) -- (-3.0, 5.45)
    node[midway, right=3pt, maplabel=arrowclr]
    {instantiate};
\node[font=\scriptsize, text=arrowclr, anchor=west] at (-2.8, 5.95)
    {$\Psi_{x_i} \in \bigoplus_k W_k$};

% Layer 2 -> Layer 3 (right side)
\draw[maparrow=arrowclr, dashed] (3.0, 4.15) -- (3.0, 3.0)
    node[midway, right=3pt, maplabel=arrowclr]
    {compute};
\node[font=\scriptsize, text=arrowclr, anchor=west] at (3.2, 3.55)
    {attention heads};

% Layer 1 -> Layer 3 (curved arrow, right side, bypassing layer 2)
\draw[maparrow=arrowclr, dotted, line width=0.8pt]
    (5.6, 6.55) to[out=-70, in=70]
    node[right=2pt, midway, font=\scriptsize, text=arrowclr]
    {direct instantiation}
    (5.6, 3.0);

% ====================================================================
% BRACE: "Standpoint Bundle" annotation on the left
% ====================================================================
\draw[decorate, decoration={brace, amplitude=6pt, mirror}, line width=0.7pt, arrowclr]
    (-6.6, 2.8) -- (-6.6, 6.2)
    node[midway, left=8pt, font=\footnotesize\bfseries, text=arrowclr, rotate=90, anchor=south]
    {Standpoint Bundle};

% ====================================================================
% RIGHT ANNOTATION: Key identity
% ====================================================================
\node[draw=arrowclr, fill=white, rounded corners=2pt,
      inner sep=5pt, font=\footnotesize, align=center, text width=3.4cm]
    at (7.2, 4.8)
    {\textbf{Key Identity:}\\[2pt]
     Attention $=$ Parallel\\Transport on $(E, X, \pi)$};

\end{tikzpicture}
\caption{%
    \textbf{LCESA framework overview.}
    The abstract geometric framework (top) defines the event space~$\mathcal{X}$, fiber~$E_x$,
    connection~$A_\Psi$, and curvature~$F$ on the standpoint bundle.
    The five-layer standpoint model (middle) instantiates the fiber as a direct sum of
    psychologically grounded subspaces~$\psi_k$.
    The transformer instantiation (bottom) maps each geometric object to a concrete
    architectural component: events to turns, fibers to head value subspaces~$W_k$,
    connection to attention weights~$\alpha$, and curvature to path-dependence of
    composite transport.
    Dashed arrows indicate the instantiation mapping between layers.
}
\label{fig:lcesa-framework}
\end{figure}

% Requires: \usepackage{tikz}
%           \usetikzlibrary{arrows.meta, positioning, shapes.geometric,
%                           decorations.pathreplacing, calc, fit}
% ==============================================================================

\begin{figure}[t]
\centering
\begin{tikzpicture}[
    % --- global styles ---
    >=Stealth,
    node distance=0.6cm and 0.4cm,
    every node/.style={font=\small},
    % --- layer box styles ---
    layerbox/.style={
        draw=#1, fill=#1!6, rounded corners=4pt,
        inner sep=8pt, minimum width=12.8cm
    },
    % --- component box styles ---
    abstractbox/.style={
        draw=abstractclr, fill=abstractclr!10, rounded corners=3pt,
        inner sep=5pt, minimum height=0.9cm, font=\small
    },
    layercomp/.style={
        draw=layerclr, fill=layerclr!10, rounded corners=3pt,
        inner sep=4pt, minimum height=0.85cm, minimum width=1.9cm,
        font=\small, align=center
    },
    instbox/.style={
        draw=instclr, fill=instclr!10, rounded corners=3pt,
        inner sep=4pt, minimum height=0.85cm, minimum width=2.1cm,
        font=\small, align=center
    },
    % --- arrow styles ---
    maparrow/.style={
        ->, line width=0.9pt, draw=#1, >=Stealth
    },
    maplabel/.style={
        font=\footnotesize, text=#1, fill=white, inner sep=1.5pt
    },
]

% === Colors ===
\definecolor{abstractclr}{RGB}{41, 98, 165}   % blue for abstract geometry
\definecolor{layerclr}{RGB}{34, 139, 87}      % green for five-layer model
\definecolor{instclr}{RGB}{180, 95, 25}       % orange for transformer instantiation
\definecolor{arrowclr}{RGB}{100, 100, 100}

% ====================================================================
% LAYER 1: Abstract Geometric Framework
% ====================================================================
\node[layerbox=abstractclr, label={[font=\small\bfseries, abstractclr]above:Abstract Geometric Framework}]
    (layer1) at (0, 7.4) {};

% Components inside layer 1
\node[abstractbox, minimum width=2.4cm] (eventspace)
    at (-4.6, 7.4) {$\mathcal{X} = (V, E, T)$};
\node[abstractbox, minimum width=1.6cm] (fiber)
    at (-1.6, 7.4) {$E_x$};
\node[abstractbox, minimum width=2.0cm] (connection)
    at (1.1, 7.4) {$A_{\Psi}$};
\node[abstractbox, minimum width=1.6cm] (curvature)
    at (3.6, 7.4) {$F$};
\node[abstractbox, minimum width=2.0cm] (gauge)
    at (5.8, 7.4) {$G = \prod_k O(d_k)$};

% Sub-labels beneath each component
\node[font=\scriptsize, text=abstractclr] at (-4.6, 6.85) {event space};
\node[font=\scriptsize, text=abstractclr] at (-1.6, 6.85) {fiber};
\node[font=\scriptsize, text=abstractclr] at (1.1, 6.85) {connection};
\node[font=\scriptsize, text=abstractclr] at (3.6, 6.85) {curvature};
\node[font=\scriptsize, text=abstractclr] at (5.8, 6.85) {gauge group};

% Arrows inside layer 1 (horizontal flow)
\draw[->, abstractclr, line width=0.6pt] (eventspace.east) -- (fiber.west);
\draw[->, abstractclr, line width=0.6pt] (fiber.east) -- (connection.west);
\draw[->, abstractclr, line width=0.6pt] (connection.east) -- (curvature.west);
\draw[->, abstractclr, line width=0.6pt] (curvature.east) -- (gauge.west);

% ====================================================================
% LAYER 2: Five-Layer Standpoint Model
% ====================================================================
\node[layerbox=layerclr, label={[font=\small\bfseries, layerclr]above:Five-Layer Standpoint Model $\Psi_t$}]
    (layer2) at (0, 4.8) {};

% Five standpoint layer components
\node[layercomp, minimum width=2.1cm] (psi_min)
    at (-4.7, 4.8) {$\psi_{\min}$\\[-2pt]{\scriptsize minimal self}};
\node[layercomp, minimum width=2.1cm] (psi_nar)
    at (-2.3, 4.8) {$\psi_{\mathrm{nar}}$\\[-2pt]{\scriptsize narrative}};
\node[layercomp, minimum width=2.1cm] (psi_soc)
    at (0.1, 4.8) {$\psi_{\mathrm{soc}}$\\[-2pt]{\scriptsize social}};
\node[layercomp, minimum width=2.1cm] (psi_mor)
    at (2.5, 4.8) {$\psi_{\mathrm{mor}}$\\[-2pt]{\scriptsize moral}};
\node[layercomp, minimum width=2.1cm] (psi_pos)
    at (4.9, 4.8) {$\psi_{\mathrm{pos}}$\\[-2pt]{\scriptsize positional}};

% Orthogonal decomposition label
\node[font=\scriptsize, text=layerclr] at (0, 4.2)
    {$\Psi_t = \psi_{\mathrm{min}} \oplus \psi_{\mathrm{nar}} \oplus \psi_{\mathrm{soc}} \oplus \psi_{\mathrm{mor}} \oplus \psi_{\mathrm{pos}}$};

% ====================================================================
% LAYER 3: Transformer Instantiation
% ====================================================================
\node[layerbox=instclr, label={[font=\small\bfseries, instclr]above:Transformer Instantiation}]
    (layer3) at (0, 2.0) {};

% Components inside layer 3
\node[instbox, minimum width=2.6cm] (turns)
    at (-4.3, 2.0) {Turns $x_i$\\[-2pt]{\scriptsize (events)}};
\node[instbox, minimum width=2.6cm] (subspaces)
    at (-1.3, 2.0) {Head subspaces\\[-2pt]{\scriptsize $W_k = \operatorname{span}(V_h)$}};
\node[instbox, minimum width=2.6cm] (attn)
    at (1.7, 2.0) {Attention $\alpha$\\[-2pt]{\scriptsize (transport $U_{ij}$)}};
\node[instbox, minimum width=2.6cm] (pathdep)
    at (4.7, 2.0) {Path-dependence\\[-2pt]{\scriptsize (curvature $F_{ijk}$)}};

% Sub-labels
\node[font=\scriptsize, text=instclr] at (-4.3, 1.35) {$V = \{x_1, \ldots, x_n\}$};
\node[font=\scriptsize, text=instclr] at (-1.3, 1.35) {$\mathbb{R}^d = \bigoplus_k W_k$};
\node[font=\scriptsize, text=instclr] at (1.7, 1.35) {$U_{ij} = \sum_k \bar\alpha_{ji}^{(k)} P_{W_k}$};
\node[font=\scriptsize, text=instclr] at (4.7, 1.35) {$F_{ijk} = U_{ij} U_{jk} (U_{ik}^{\exp})^{-1}$};

% Arrows inside layer 3
\draw[->, instclr, line width=0.6pt] (turns.east) -- (subspaces.west);
\draw[->, instclr, line width=0.6pt] (subspaces.east) -- (attn.west);
\draw[->, instclr, line width=0.6pt] (attn.east) -- (pathdep.west);

% ====================================================================
% INTER-LAYER MAPPING ARROWS
% ====================================================================

% Layer 1 -> Layer 2 (left side)
\draw[maparrow=arrowclr, dashed] (-3.0, 6.55) -- (-3.0, 5.45)
    node[midway, right=3pt, maplabel=arrowclr]
    {instantiate};
\node[font=\scriptsize, text=arrowclr, anchor=west] at (-2.8, 5.95)
    {$\Psi_{x_i} \in \bigoplus_k W_k$};

% Layer 2 -> Layer 3 (right side)
\draw[maparrow=arrowclr, dashed] (3.0, 4.15) -- (3.0, 3.0)
    node[midway, right=3pt, maplabel=arrowclr]
    {compute};
\node[font=\scriptsize, text=arrowclr, anchor=west] at (3.2, 3.55)
    {attention heads};

% Layer 1 -> Layer 3 (curved arrow, right side, bypassing layer 2)
\draw[maparrow=arrowclr, dotted, line width=0.8pt]
    (5.6, 6.55) to[out=-70, in=70]
    node[right=2pt, midway, font=\scriptsize, text=arrowclr]
    {direct instantiation}
    (5.6, 3.0);

% ====================================================================
% BRACE: "Standpoint Bundle" annotation on the left
% ====================================================================
\draw[decorate, decoration={brace, amplitude=6pt, mirror}, line width=0.7pt, arrowclr]
    (-6.6, 2.8) -- (-6.6, 6.2)
    node[midway, left=8pt, font=\footnotesize\bfseries, text=arrowclr, rotate=90, anchor=south]
    {Standpoint Bundle};

% ====================================================================
% RIGHT ANNOTATION: Key identity
% ====================================================================
\node[draw=arrowclr, fill=white, rounded corners=2pt,
      inner sep=5pt, font=\footnotesize, align=center, text width=3.4cm]
    at (7.2, 4.8)
    {\textbf{Key Identity:}\\[2pt]
     Attention $=$ Parallel\\Transport on $(E, X, \pi)$};

\end{tikzpicture}
\caption{%
    \textbf{LCESA framework overview.}
    The abstract geometric framework (top) defines the event space~$\mathcal{X}$, fiber~$E_x$,
    connection~$A_\Psi$, and curvature~$F$ on the standpoint bundle.
    The five-layer standpoint model (middle) instantiates the fiber as a direct sum of
    psychologically grounded subspaces~$\psi_k$.
    The transformer instantiation (bottom) maps each geometric object to a concrete
    architectural component: events to turns, fibers to head value subspaces~$W_k$,
    connection to attention weights~$\alpha$, and curvature to path-dependence of
    composite transport.
    Dashed arrows indicate the instantiation mapping between layers.
}
\label{fig:lcesa-framework}
\end{figure}

% Requires: \usepackage{tikz}
%           \usetikzlibrary{arrows.meta, positioning, shapes.geometric,
%                           decorations.pathreplacing, calc, fit}
% ==============================================================================

\begin{figure}[t]
\centering
\begin{tikzpicture}[
    % --- global styles ---
    >=Stealth,
    node distance=0.6cm and 0.4cm,
    every node/.style={font=\small},
    % --- layer box styles ---
    layerbox/.style={
        draw=#1, fill=#1!6, rounded corners=4pt,
        inner sep=8pt, minimum width=12.8cm
    },
    % --- component box styles ---
    abstractbox/.style={
        draw=abstractclr, fill=abstractclr!10, rounded corners=3pt,
        inner sep=5pt, minimum height=0.9cm, font=\small
    },
    layercomp/.style={
        draw=layerclr, fill=layerclr!10, rounded corners=3pt,
        inner sep=4pt, minimum height=0.85cm, minimum width=1.9cm,
        font=\small, align=center
    },
    instbox/.style={
        draw=instclr, fill=instclr!10, rounded corners=3pt,
        inner sep=4pt, minimum height=0.85cm, minimum width=2.1cm,
        font=\small, align=center
    },
    % --- arrow styles ---
    maparrow/.style={
        ->, line width=0.9pt, draw=#1, >=Stealth
    },
    maplabel/.style={
        font=\footnotesize, text=#1, fill=white, inner sep=1.5pt
    },
]

% === Colors ===
\definecolor{abstractclr}{RGB}{41, 98, 165}   % blue for abstract geometry
\definecolor{layerclr}{RGB}{34, 139, 87}      % green for five-layer model
\definecolor{instclr}{RGB}{180, 95, 25}       % orange for transformer instantiation
\definecolor{arrowclr}{RGB}{100, 100, 100}

% ====================================================================
% LAYER 1: Abstract Geometric Framework
% ====================================================================
\node[layerbox=abstractclr, label={[font=\small\bfseries, abstractclr]above:Abstract Geometric Framework}]
    (layer1) at (0, 7.4) {};

% Components inside layer 1
\node[abstractbox, minimum width=2.4cm] (eventspace)
    at (-4.6, 7.4) {$\mathcal{X} = (V, E, T)$};
\node[abstractbox, minimum width=1.6cm] (fiber)
    at (-1.6, 7.4) {$E_x$};
\node[abstractbox, minimum width=2.0cm] (connection)
    at (1.1, 7.4) {$A_{\Psi}$};
\node[abstractbox, minimum width=1.6cm] (curvature)
    at (3.6, 7.4) {$F$};
\node[abstractbox, minimum width=2.0cm] (gauge)
    at (5.8, 7.4) {$G = \prod_k O(d_k)$};

% Sub-labels beneath each component
\node[font=\scriptsize, text=abstractclr] at (-4.6, 6.85) {event space};
\node[font=\scriptsize, text=abstractclr] at (-1.6, 6.85) {fiber};
\node[font=\scriptsize, text=abstractclr] at (1.1, 6.85) {connection};
\node[font=\scriptsize, text=abstractclr] at (3.6, 6.85) {curvature};
\node[font=\scriptsize, text=abstractclr] at (5.8, 6.85) {gauge group};

% Arrows inside layer 1 (horizontal flow)
\draw[->, abstractclr, line width=0.6pt] (eventspace.east) -- (fiber.west);
\draw[->, abstractclr, line width=0.6pt] (fiber.east) -- (connection.west);
\draw[->, abstractclr, line width=0.6pt] (connection.east) -- (curvature.west);
\draw[->, abstractclr, line width=0.6pt] (curvature.east) -- (gauge.west);

% ====================================================================
% LAYER 2: Five-Layer Standpoint Model
% ====================================================================
\node[layerbox=layerclr, label={[font=\small\bfseries, layerclr]above:Five-Layer Standpoint Model $\Psi_t$}]
    (layer2) at (0, 4.8) {};

% Five standpoint layer components
\node[layercomp, minimum width=2.1cm] (psi_min)
    at (-4.7, 4.8) {$\psi_{\min}$\\[-2pt]{\scriptsize minimal self}};
\node[layercomp, minimum width=2.1cm] (psi_nar)
    at (-2.3, 4.8) {$\psi_{\mathrm{nar}}$\\[-2pt]{\scriptsize narrative}};
\node[layercomp, minimum width=2.1cm] (psi_soc)
    at (0.1, 4.8) {$\psi_{\mathrm{soc}}$\\[-2pt]{\scriptsize social}};
\node[layercomp, minimum width=2.1cm] (psi_mor)
    at (2.5, 4.8) {$\psi_{\mathrm{mor}}$\\[-2pt]{\scriptsize moral}};
\node[layercomp, minimum width=2.1cm] (psi_pos)
    at (4.9, 4.8) {$\psi_{\mathrm{pos}}$\\[-2pt]{\scriptsize positional}};

% Orthogonal decomposition label
\node[font=\scriptsize, text=layerclr] at (0, 4.2)
    {$\Psi_t = \psi_{\mathrm{min}} \oplus \psi_{\mathrm{nar}} \oplus \psi_{\mathrm{soc}} \oplus \psi_{\mathrm{mor}} \oplus \psi_{\mathrm{pos}}$};

% ====================================================================
% LAYER 3: Transformer Instantiation
% ====================================================================
\node[layerbox=instclr, label={[font=\small\bfseries, instclr]above:Transformer Instantiation}]
    (layer3) at (0, 2.0) {};

% Components inside layer 3
\node[instbox, minimum width=2.6cm] (turns)
    at (-4.3, 2.0) {Turns $x_i$\\[-2pt]{\scriptsize (events)}};
\node[instbox, minimum width=2.6cm] (subspaces)
    at (-1.3, 2.0) {Head subspaces\\[-2pt]{\scriptsize $W_k = \operatorname{span}(V_h)$}};
\node[instbox, minimum width=2.6cm] (attn)
    at (1.7, 2.0) {Attention $\alpha$\\[-2pt]{\scriptsize (transport $U_{ij}$)}};
\node[instbox, minimum width=2.6cm] (pathdep)
    at (4.7, 2.0) {Path-dependence\\[-2pt]{\scriptsize (curvature $F_{ijk}$)}};

% Sub-labels
\node[font=\scriptsize, text=instclr] at (-4.3, 1.35) {$V = \{x_1, \ldots, x_n\}$};
\node[font=\scriptsize, text=instclr] at (-1.3, 1.35) {$\mathbb{R}^d = \bigoplus_k W_k$};
\node[font=\scriptsize, text=instclr] at (1.7, 1.35) {$U_{ij} = \sum_k \bar\alpha_{ji}^{(k)} P_{W_k}$};
\node[font=\scriptsize, text=instclr] at (4.7, 1.35) {$F_{ijk} = U_{ij} U_{jk} (U_{ik}^{\exp})^{-1}$};

% Arrows inside layer 3
\draw[->, instclr, line width=0.6pt] (turns.east) -- (subspaces.west);
\draw[->, instclr, line width=0.6pt] (subspaces.east) -- (attn.west);
\draw[->, instclr, line width=0.6pt] (attn.east) -- (pathdep.west);

% ====================================================================
% INTER-LAYER MAPPING ARROWS
% ====================================================================

% Layer 1 -> Layer 2 (left side)
\draw[maparrow=arrowclr, dashed] (-3.0, 6.55) -- (-3.0, 5.45)
    node[midway, right=3pt, maplabel=arrowclr]
    {instantiate};
\node[font=\scriptsize, text=arrowclr, anchor=west] at (-2.8, 5.95)
    {$\Psi_{x_i} \in \bigoplus_k W_k$};

% Layer 2 -> Layer 3 (right side)
\draw[maparrow=arrowclr, dashed] (3.0, 4.15) -- (3.0, 3.0)
    node[midway, right=3pt, maplabel=arrowclr]
    {compute};
\node[font=\scriptsize, text=arrowclr, anchor=west] at (3.2, 3.55)
    {attention heads};

% Layer 1 -> Layer 3 (curved arrow, right side, bypassing layer 2)
\draw[maparrow=arrowclr, dotted, line width=0.8pt]
    (5.6, 6.55) to[out=-70, in=70]
    node[right=2pt, midway, font=\scriptsize, text=arrowclr]
    {direct instantiation}
    (5.6, 3.0);

% ====================================================================
% BRACE: "Standpoint Bundle" annotation on the left
% ====================================================================
\draw[decorate, decoration={brace, amplitude=6pt, mirror}, line width=0.7pt, arrowclr]
    (-6.6, 2.8) -- (-6.6, 6.2)
    node[midway, left=8pt, font=\footnotesize\bfseries, text=arrowclr, rotate=90, anchor=south]
    {Standpoint Bundle};

% ====================================================================
% RIGHT ANNOTATION: Key identity
% ====================================================================
\node[draw=arrowclr, fill=white, rounded corners=2pt,
      inner sep=5pt, font=\footnotesize, align=center, text width=3.4cm]
    at (7.2, 4.8)
    {\textbf{Key Identity:}\\[2pt]
     Attention $=$ Parallel\\Transport on $(E, X, \pi)$};

\end{tikzpicture}
\caption{%
    \textbf{LCESA framework overview.}
    The abstract geometric framework (top) defines the event space~$\mathcal{X}$, fiber~$E_x$,
    connection~$A_\Psi$, and curvature~$F$ on the standpoint bundle.
    The five-layer standpoint model (middle) instantiates the fiber as a direct sum of
    psychologically grounded subspaces~$\psi_k$.
    The transformer instantiation (bottom) maps each geometric object to a concrete
    architectural component: events to turns, fibers to head value subspaces~$W_k$,
    connection to attention weights~$\alpha$, and curvature to path-dependence of
    composite transport.
    Dashed arrows indicate the instantiation mapping between layers.
}
\label{fig:lcesa-framework}
\end{figure}
